\section{Explicit domain, assignment and closure proofs for the supporting comparison}

\begin{annexitem}{annex:p8.kc.evaluation}{thm:p08-kc-evaluation}{Truth recursion with exactly the supplied free-variable assignment}
\begin{proof}
Induct on formula complexity, simultaneously for every world and
every assignment of the formula's free variables into U. Atomic
truth is uniquely determined by the assigned tuple's membership
in its given U-extension; membership is specified whether or
not the tuple's members exist at that world. The zero-place
case has its independently specified T/F value. Classical
negation and conjunction uniquely determine their result from
the inner results. For box, all accessible worlds use the same
assignment of the same free variables, so the already determined
truth values say uniquely whether all are T. For \((x)A\),
the supplied assignment of the other free variables is extended
by each \(a\in\psi(H)\), a subset of U. The induction then
gives each inner truth value; the universal condition determines
the quantified value. If that domain is empty, the universal
condition is true. No global assignment to infinitely many
variables is needed: for a closed formula its empty assignment
exists even when U is empty.

The same induction proves dependence only on free-variable
values. In the quantifier step two agreeing assignments of the
remaining variables yield agreeing extensions for each a. For
capture-free variable replacement, compare the substituted
formula's assignment with the original after updating x to
the assigned value of y. At atoms the tuples are equal. Boolean
operations and the box case preserve equality, since variable
values do not change with worlds. At a quantifier binding z,
rename z if it would capture a newly substituted y; its updates
then commute with the x-update and the inner induction applies
to each local a. If the quantifier binds x, no free x under it
is replaced and assignment irrelevance supplies the result.
These cases prove the stated replacement lemma. Neither the
modal nor atomic step assumed current-domain membership of
an unchanged free-variable value.
\end{proof}
\end{annexitem}

\begin{annexitem}{annex:p8.kc.closed-soundness}{thm:p08-kc-closed-soundness}{Every closure schema and closed derived necessitation}
\begin{proof}
Call a formula fully true when it is true at every world for
every assignment of its free variables into U. If U is empty
and the formula has free variables, this quantification over
assignments is empty; its subsequent quantified closures are
handled by the same explicit local-domain clauses.
Truth-functional tautologies are fully true because every
supplied subformula evaluation is T or F. Reflexivity makes
\(\Box A\to A\) fully true. For distribution, an assignment
making \(\Box(A\to B)\) and \(\Box A\) true gives at
each accessible world both the implication and its antecedent,
so B is true at each and \(\Box B\) follows. If x is not
free in A, truth of A is independent of any x-update, proving
\(A\to(x)A\), including the empty-domain case. Quantifier
distribution follows pointwise: if all current a satisfy
\(A\to B\) and all satisfy A, all satisfy B.

These first five matrices are fully true, so any prefix of
universal quantifiers and boxes preserves their full truth.
For universal quantification, each updated value is in U;
for box, every successor evaluates the same supplied assignment.
The last source matrix is already \((y)((x)A(x)\to A(y))\).
For each local y=a, truth of \((x)A(x)\) gives truth at that
same a, and variable replacement identifies the A(y) instance.
Thus this whole quantified matrix is fully true for any other
free parameters. The \emph{unquantified} instantiation matrix
need not be fully true for y assigned outside the local domain;
no such claim was used. Prefixing further quantifiers and boxes
therefore gives all its permitted closures. The source extensions'
S4 matrix is fully true on transitive frames: if \(\Box A\)
holds at H and \(HRH'RH''\), transitivity makes \(HRH''\),
so A holds at \(H''\); each \(H'\) satisfies \(\Box A\)
and H satisfies its repeated box. For the Brouwersche matrix,
if A holds at H and \(HRH'\), symmetry makes H accessible
from \(H'\), so \(\Diamond A\) holds there and its box
holds at H. For S5, if \(\Diamond A\) holds at H, choose
an accessible J satisfying A. For any \(HRH'\), symmetry
and transitivity give \(H'RJ\), so \(H'\) satisfies the
same diamond; hence H satisfies \(\Box\Diamond A\).
Each argument keeps the free assignment unchanged, so the
same closure argument applies to these matrices.
Detachment of two valid closed formulas preserves truth
at the distinguished world of every model. Since any world may
be distinguished without changing the other structure, this
theorem validity also holds at all worlds of each model.

For admissibility of necessitation, induct over a finite closed
proof. If a line is a closure axiom, prefixing one more box
leaves a closure axiom. If a line B follows by detachment from
closed A and \(A\to B\), induction gives the boxed versions
of those two theorems. The closed distribution axiom and two
detachments give \(\Box B\). This is a finite theorem proof
transformation, not necessitation of a line dependent on local
premises. A closed theorem also gives \((x)A\) using the
closed instance of \(A\to(x)A\); all variables are absent
from a closed A. A free-formula assertion must first be read
as its specified closed formula before either derived operation
is applied.
\end{proof}
\end{annexitem}

\begin{annexitem}{annex:p8.kc.necessitation-boundary}{thm:p08-kc-necessitation-boundary}{Two different closures and the circular use of converse Barcan}
\begin{proof}
Use the second full S5 model in the body: worlds G,H, complete
R, \(\psi(G)=\{a,b\}\), \(\psi(H)=\{a\}\),
\(P_G=\{a,b\}\), \(P_H=\{a\}\). Put
\(C(y)=((x)P(x)\to P(y))\), with x,y distinct.
At G every locally assigned y is in \(\{a,b\}\), and both
satisfy P at G. At H its sole locally assigned y=a satisfies
P there. Hence \((y)C(y)\) is true at both worlds and
\(\Box(y)C(y)\) is true at G. But at H, under the unchanged
assignment y=b from G, \((x)P(x)\) is true and P(b) false.
Thus C(b) fails at H, \(\Box C(b)\) fails at G, and
\((y)\Box C(y)\) fails at G. This directly separates the
two closures in the actual source example.

The seemingly classical converse-Barcan deduction begins with
the instantiation assertion \((x)A(x)\to A(y)\), necessitates
it, applies box distribution, and generalizes over y. Under
the source generality interpretation, its first line denotes
\((y)((x)A(x)\to A(y))\). Necessitation therefore supplies
\(\Box(y)((x)A(x)\to A(y))\). The later distribution
step would need \((y)\Box((x)A(x)\to A(y))\), which is
not that supplied theorem. Converting between them is exactly
a CBF instance for C(y). The displayed model falsifies that
instance. Therefore the deduction is circular at this change
of closure order, before its sound subsequent distribution
or generalization steps can contribute.
\end{proof}
\end{annexitem}

\begin{annexitem}{annex:p8.kc.barcan-domains}{thm:p08-kc-barcan-domains}{Both directions, empty cases and exact arbitrary-interpretation converses}
\begin{proof}
Fix any other free-variable assignment s at H. Assume domains
decrease along edges and \((x)\Box A(x)\) is true at H.
For any \(H'\) with \(HRH'\), every
\(a\in\psi(H')\) is also in \(\psi(H)\). The antecedent
therefore made \(\Box A\) true at H under s[x\(\mapsto\)a],
and in particular A is true at \(H'\) under that same
assignment. All its local a satisfy A, so \((x)A\) is true
at \(H'\); this proves \(\Box(x)A\) and BF. If any
domain is empty, the same universal assertions have their
declared vacuous meanings; no witness outside U is invented.

For CBF assume domains increase and \(\Box(x)A\) is true
at H. For any \(a\in\psi(H)\) and any accessible \(H'\),
the same a belongs to \(\psi(H')\), so the boxed quantified
premise makes A true there under s[x\(\mapsto\)a]. Thus
\(\Box A\) is true at H for every current a and
\((x)\Box A\) follows. This proves CBF, again including
empty domains. Constants are not needed in either direction.

For the BF converse, suppose an edge \(HRH'\) has
\(b\in\psi(H')\setminus\psi(H)\). Set the unary
predicate extension at \(H'\) to U\(\setminus\)\{b\},
and at every other world to U. Every root-domain a satisfies
P at every root successor, since a\(\ne\)b; hence
\((x)\Box P(x)\) is true at H. At \(H'\), its existing
b falsifies P and \((x)P(x)\); BF is false at H. A failure
of inclusion cannot be a reflexive self-edge, so no conflict
between the two extension prescriptions occurs.

For the CBF converse, suppose an edge has
\(b\in\psi(H)\setminus\psi(H')\). Set the unary
predicate at each world J to \(\psi(J)\), as a subset of U.
Every local universal \((x)P(x)\) is true, so its box is
true at H. At \(H'\), P(b) is false, so \(\Box P(b)\)
is false at H and the root quantifier, which includes b,
falsifies \((x)\Box P(x)\). These valuations prove that
all-interpretation BF/CBF validity forces exactly their stated
edge inclusions. Combining the implications gives the full
table: both inclusions ensure both directions; each isolated
inclusion ensures only its corresponding one; and either
failed inclusion has an actual unary counterinterpretation.
The proofs concern varying-domain truth with unchanged
U-valued variables, not a convention transporting assignments
to new objects at successive worlds.
\end{proof}
\end{annexitem}

\begin{annexitem}{annex:p8.kc.s5-countermodels}{thm:p08-kc-s5-countermodels}{Every relation pair, local quantifier and absent-object evaluation}
\begin{proof}
In the first model R consists of the four pairs
\((G,G),(G,H),(H,G),(H,H)\), so is an equivalence relation.
The union U is \(\{a,b\}\). The predicate is true of a
at both worlds and false of b at both. At G the only local
quantifier value is a, and \(\Box P(a)\) is true because
both successors satisfy it. Thus the BF antecedent is true.
At H its local b falsifies P, so \((x)P(x)\) fails there
and the BF consequent's box fails at G. This computation
never needed P(b) at G; the failure survives alternative
policies which leave that one absent-object assertion unevaluated,
provided the remaining displayed evaluations and modal choices
remain as stated.

In the second model the same four relation pairs and the same
U apply, but P is true of both a,b at G and only a at H.
Every locally existing object satisfies P in each world, so
\(\Box(x)P(x)\) is true at G. The local b at G is
nevertheless a false P-argument at H, under the chosen total
U-extension convention. Hence \(\Box P(b)\) is false at G
and its current universal closure fails, refuting CBF. This
failure does depend on evaluating the absent-object P(b) at H
as false. A truth-gap semantics with necessity meaning merely
``never false'' may change that result; it is not the chosen
two-valued clause. Finally the reverse edge in each example
has the domain inclusion opposite to its forward edge. They
are S5 countermodels to unrestricted domains, not exceptions
to the proved globally increasing/decreasing classifications.
\end{proof}
\end{annexitem}

\begin{annexitem}{annex:p8.kc.existence}{thm:p08-kc-existence}{Existence guards, identity and the explicit strict-predicate substitution failure}
\begin{proof}
At H, every \(a\in\psi(H)\) is in E's extension, so
\((x)E(x)\) is true. If \((x)A(x)\) is true and E(y)
holds under a supplied assignment, the assigned y belongs
to \(\psi(H)\), and its A(y) is true. This makes the
guarded instantiation matrix fully true; prefixes of boxes
and quantifiers give its closures by the soundness argument.
For identity, \((\exists y)(x=y)\) is true iff some local
\(a\in\psi(H)\) equals the unchanged U-value assigned
to x, exactly iff E(x) is true. This includes empty domains
and absent assigned values.

Every world's local universal E is true, so \(\Box(x)E(x)\)
is valid. In the second S5 countermodel set E to its required
domain extension, which is precisely that model's P. Its
\((x)\Box E(x)\) is false at G for the same b witness.
Classical excluded middle \(A(x)\lor\neg A(x)\) is true
for every supplied U-assignment, whether or not x exists there.
It is therefore not extensionally identical to E(x).

For the strict-predicate alternative, consider the closed formula
\[
 (x)\Box\bigl(P(x)\land(y)Q(y)\to Q(x)\bigr).
\]
If P and Q have extensions contained in each current domain,
the inner implication is true at every world for every U-value
of x: if P(x) holds, that value is local; the local universal
Q then entails Q(x). Therefore its box and root quantifier
are true on all such strict-predicate models.
But substitute the formula \(Q(x)\lor\neg Q(x)\) for P(x).
Take worlds G,H with reflexive loops and edge \(GRH\),
\(\psi(G)=\{a,b\},\psi(H)=\{a\}\), and
\(Q_G=\{a,b\},Q_H=\{a\}\). At H the substituted
antecedent is true under x=b, because the substituted P
formula is a tautology and its local universal Q is true.
Q(b) is false there. Thus the substituted boxed matrix fails
at G under x=b, and its root universal is false. The remaining
atomic Q still satisfies the strict extension policy. This
is an actual failure of unrestricted formula substitution in
that alternative class, not a failure in the source's broader
arbitrary-U-extension class, where the unsubstituted schema
was never generally valid.
\end{proof}
\end{annexitem}

\begin{annexitem}{annex:p8.kc.fitting-bridge}{thm:p08-kc-fitting-bridge}{All translation cases and the separate fixed-domain calculation}
\begin{proof}
Use the same worlds and object union, set the Fitting world
domains to \(\psi(H)\), and preserve each predicate extension.
Set its E extension to that same domain. Given an assignment
of a formula's free variables, extend it to other variables
only if needed in the consumer's valuation notation; since
the nonempty local-domain hypothesis makes U nonempty, a
default object supplies such an extension, and assignment
irrelevance ensures that it changes no truth value.

Induct on source formulas. At atoms the same assigned tuple
has the same membership in the same extension. Boolean
clauses agree. For box the worlds reached and the unchanged
assignment values agree, so the inner truth induction agrees
at every successor. For source \((x)A\), the translation
\((\forall x)(E(x)\to t(A))\) uses Fitting's primitive
quantifier over U. Its E guard is true exactly for
\(a\in\psi(H)\); all other U-values make the implication
true. For the guarded values the inner induction gives exactly
the source truth test. This proves the local universal step,
and existential follows by its Boolean expansion. Thus the
entire common-language translation preserves truth. A source
identity atom, if the optional matching extension is selected,
has the same two assigned U-values in both models and hence the
same equality truth value. Without that matching clause the
identity extension is not part of the claimed common language.
A source
model with an empty local domain does not satisfy that consumer's
declared nonempty-domain requirement; broadening the consumer
would be a different comparison, not something presumed here.

For fixed-U primitive quantification, evaluate
\(\forall x\Box A(x)\) at H. It says A is true for
every a in U and every accessible successor, with each same
updated assignment. \(\Box\forall xA(x)\) states precisely
the same pair of universal conditions in the opposite order.
They are equivalent because the quantifier range is fixed U,
so both implication directions hold, independently of how
the separate \(\psi\) and E vary. This is not the source
local-domain binder and not the actualist table. The formal
definitions establish only these semantic dependencies, not
which quantified or metaphysical position is philosophically correct.
\end{proof}
\end{annexitem}

\begin{annexitem}{annex:p8.kc.designation}{thm:p08-kc-designation}{Rigidity without necessary existence and a nonrigid substitution witness}
\begin{proof}
For the explicitly added constant semantics, let
\(K=\{G,H\}\), R have both reflexive loops and edge
\(GRH\), \(\psi(G)=\{a,b\},\psi(H)=\{b\}\),
and let c rigidly denote a. Then E(c) is true at G, false
at H, and \(\Box E(c)\) false at G. At both worlds
c=c is true, because identity compares denotation in U rather
than the current existence domain; its box is true. This
separates fixed designation from necessary existence without
changing the source's variable clauses. A name taking two
different values, each within its respective local domain,
would conversely exist everywhere while failing rigidity.

For the substitution warning use constant domains
\(\psi(G)=\psi(H)=U=\{a,b\}\), the same reflexive
relation with edge G to H, and d rigidly denoting a. Let
c denote a at G and b at H. For source-style fixed variables,
\(x=y\to\Box(x=y)\) is true under every U-assignment:
if the two values agree they continue to agree at each successor;
otherwise the implication has false antecedent. Its ordinary
universal closure is valid as well. But the named instance
\(c=d\to\Box(c=d)\) is false at G: its current equality
is true and its equality at H false. Thus variable rigidity
does not prove a theorem-preserving substitution principle for
nonrigid names. Both Barcan directions still hold in this
constant-domain model by the separate domain argument. This
is explicitly an added designation-language countermodel,
not a claimed variable-only 1963 source counterexample.
\end{proof}
\end{annexitem}
